%! Characteristics of Polynomial Functions — Unit 4, Lesson 4.0 — Exit Ticket (Answer Key)
\documentclass[10pt]{article}
\usepackage{algebra2-article}
\usepackage{algebra2-key}   % pulls in -boxes; do NOT also load -boxes
\newcommand{\TallMath}[1]{$\displaystyle #1\rule[-1.4em]{0pt}{3.2em}$}
\newcommand{\lgrid}[4]{%
  \draw[step=1, linegray!40, very thin] (#1,#3) grid (#2,#4);
  \draw[->, semithick, charcoal] (#1,0)--(#2,0) node[below right, font=\scriptsize]{$x$};
  \draw[->, semithick, charcoal] (0,#3)--(0,#4) node[left, font=\scriptsize]{$y$};}

\begin{document}

\pageheader{Unit 4, Lesson 4.0}{Exit Ticket --- Answer Key}
\namedateperiod

\vspace{0.10in}

No notes. The graph below is the polynomial $k(x) = (x+1)^2(x-2)$. Use it for all questions.
\begin{center}
\begin{tikzpicture}[scale=0.52]
  \lgrid{-3}{3}{-5}{3}
  \begin{scope}
    \clip (-3,-5) rectangle (3,3);
    \draw[very thick, forest, domain=-2.1:2.25, samples=90, smooth]
      plot (\x,{(\x)*(\x)*(\x) - 3*(\x) - 2});
  \end{scope}
  \fill[charcoal] (-1,0) circle (3pt) node[above left, font=\scriptsize]{$(-1,0)$};
  \fill[charcoal] (1,-4) circle (3pt) node[below right, font=\scriptsize]{$(1,-4)$};
  \fill[charcoal] (2,0) circle (3pt) node[above right, font=\scriptsize]{$(2,0)$};
  \fill[charcoal] (0,-2) circle (3pt) node[right, font=\scriptsize]{$(0,-2)$};
\end{tikzpicture}
\end{center}

\begin{enumerate}[leftmargin=1.8em, itemsep=9pt]
  \item \textbf{Read the features.} State the end behavior (both arms), the two turning points and
        whether each is a relative maximum or minimum, and the zeros --- saying whether the graph
        \emph{crosses} or \emph{touches} at each.
        \ansline{End behavior: as $x\to-\infty,\ k\to-\infty$; as $x\to+\infty,\ k\to+\infty$ (odd degree, $a>0$). Turning points: relative maximum $(-1,0)$, relative minimum $(1,-4)$. Zeros: at $x=-1$ (multiplicity $2$) the graph \emph{touches}; at $x=2$ (multiplicity $1$) it \emph{crosses}.}
  \item \textbf{Classify \& justify.} Is $k$ even, odd, or neither? Justify your answer (think about
        the degree and about how $k(-x)$ compares to $k(x)$). Does $k$ have an \emph{absolute}
        maximum? Explain.
        \ansline{Neither: $k(x)=x^3-3x-2$ mixes an odd-power term with a constant, and $k(-x)=-x^3+3x-2$ equals neither $k(x)$ nor $-k(x)$. No absolute maximum --- the right arm rises to $+\infty$, so no output is largest (the $(-1,0)$ peak is only a \emph{relative} max).}
  \item \textbf{SOL-style multiple choice.} Which statement about $k(x)=(x+1)^2(x-2)$ is \emph{true}?
        \begin{enumerate}[label=\Alph*., leftmargin=1.9em, itemsep=1pt]
          \item The graph crosses the $x$-axis at $x=-1$.
          \item $k$ has an absolute maximum.
          \item As $x\to-\infty$, $k(x)\to-\infty$. \textcolor{keyred}{\textbf{$\leftarrow$ correct}}
          \item $k$ is an even function.
        \end{enumerate}
        \begin{itemize}[leftmargin=1.4em, nosep]
          \item[] \textbf{A} is false --- the factor $(x+1)^2$ has even multiplicity, so the graph \emph{touches} at $x=-1$. \textbf{B} is false --- odd degree means no absolute max. \textbf{D} is false --- $k$ is neither even nor odd. \textbf{C} is correct: odd degree with a positive lead sends the left arm to $-\infty$.
        \end{itemize}
\end{enumerate}

\begin{teachernote}
Sort responses into ``got it / multiplicity slip (calls $x=-1$ a cross) / relative-vs-absolute slip
(calls $(-1,0)$ the maximum) / end-behavior slip.'' A wrong answer of A or D on item~3 means the
student has not connected \emph{even multiplicity $\to$ touch} or the $f(-x)$ symmetry test. If
several miss it, reopen Lesson~4.1 by re-marking the touch at $x=-1$ and the opposite arms.
\end{teachernote}

\end{document}
